\section*{Capítulo \ref{ch:g7:earthquakes} --- Los terremotos}

\begin{solution}{exo:g7:earthquakes:1}
(a) \omterm{def:g7:earthquakes:focus}{hipocentro}; (b) \omterm{def:g7:earthquakes:intensity}{intensidad}; (c) \omterm{def:g7:earthquakes:fault}{falla}; (d) \omterm{def:g7:earthquakes:focus}{epicentro}; (e) \omterm{def:g7:earthquakes:magnitude}{magnitud};
(f) \omterm{def:g7:earthquakes:seismometer}{sismómetro}.
\end{solution}

\begin{solution}{exo:g7:earthquakes:2}
(a) En el propio pueblo: el \omterm{def:g7:earthquakes:focus}{epicentro} es el punto de la superficie situado
justo encima del \omterm{def:g7:earthquakes:focus}{hipocentro}. (b) \qty{30}{km}, en vertical hacia abajo.
(c) En general, menor: la sacudida se debilita con la distancia. (También
depende del terreno: sobre \omterm{def:g3:stones-pebbles-sand:sand}{barro} blando, el segundo pueblo podría temblar
mucho de todos modos.)
\end{solution}

\begin{solution}{exo:g7:earthquakes:3}
(a) \omterm{def:g7:earthquakes:magnitude}{Magnitud} (un número con un decimal). (b) \omterm{def:g7:earthquakes:intensity}{Intensidad} (lo que se vio en
un lugar). (c) \omterm{def:g7:earthquakes:intensity}{Intensidad} (números romanos, distintos en dos lugares).
(d) \omterm{def:g7:earthquakes:magnitude}{Magnitud}: la energía liberada, un paso más de \omterm{def:g7:earthquakes:magnitude}{magnitud}.
\end{solution}

\begin{solution}{exo:g7:earthquakes:4}
(a) Falso: los científicos no pueden predecir los \omterm{def:g4:volcanoes-earthquakes-first-look:earthquake}{terremotos}; solo saben
qué \omterm{def:g7:earthquakes:fault}{fallas} volverán a romperse algún día. (b) Falso: en alta \omterm{def:g2:water-on-earth:ocean}{mar}, un
\omterm{def:g7:earthquakes:tsunami}{tsunami} rara vez mide más de un metro de altura; crece cerca de la \omterm{def:g1:landscapes-around-us:water}{costa}.
(c) Verdadero. (d) Verdadero: el 92~\% de los \omterm{def:g4:volcanoes-earthquakes-first-look:earthquake}{terremotos} de \omterm{def:g7:earthquakes:magnitude}{magnitud}~6 o
más se producen a menos de \qty{300}{km} de un \omterm{def:g7:plate-tectonics:plate-boundary}{límite de placas}.
\end{solution}

\begin{solution}{exo:g7:earthquakes:5}
(a) Unos 127 al año de \omterm{def:g7:earthquakes:magnitude}{magnitud} 6 a 6.9; unos 13 de \omterm{def:g7:earthquakes:magnitude}{magnitud} 7 a 7.9.
(b) $126.7 \div 13.12 \approx 9.7$: unas 10~veces menos. (c)
$13.12 + 0.76 + 0.04 = 13.92$: unos \textbf{14} al año. (d)
$0.04 \times 50 = 2$: dos en 50~años (Sumatra 2004 y Japón 2011).
\end{solution}

\begin{solution}{exo:g7:earthquakes:6}
(a) Un paso: 32~veces. (b) Dos pasos: $32 \times 32 = 1024$, unas
1000~veces. (c) Cada paso hace las ondas 10~veces más grandes: $10 \times 10 =
100$~veces más grandes.
\end{solution}

\begin{solution}{exo:g7:earthquakes:7}
(a) $8 \times 5 = \qty{40}{km}$. (b) $8 \times 25 = \qty{200}{km}$. (c)
$120 \div 8 = \qty{15}{s}$. (d) Una estación solo da una distancia, no una
dirección: todos los puntos situados a esa distancia forman un círculo
alrededor de la estación, y el \omterm{def:g7:earthquakes:focus}{epicentro} podría estar en cualquier punto
de él.
\end{solution}

\begin{solution}{exo:g7:earthquakes:8}
(a) V. (b) II. (c) VIII. (d) La \omterm{def:g7:earthquakes:intensity}{intensidad} cambia de un lugar a otro: el
testimonio (c) procede de un lugar cercano al \omterm{def:g7:earthquakes:focus}{epicentro}, o construido sobre
terreno blando; el testimonio (a), de un lugar más alejado, o sobre \omterm{def:g3:stones-pebbles-sand:rock}{roca}
sólida. El \omterm{def:g4:volcanoes-earthquakes-first-look:earthquake}{terremoto} sigue teniendo una sola \omterm{def:g7:earthquakes:magnitude}{magnitud}.
\end{solution}

\begin{solution}{exo:g7:earthquakes:9}
(a) $7031 - 5561 = 1470$. (b) $1470 \div 7031 \approx 0.209$: alrededor
del \textbf{21~\%}. (c) La \omterm{def:g7:plate-tectonics:plate}{placa} que se hunde baja hasta muy dentro del
\omterm{def:g6:inside-the-earth:layers}{manto}. Todavía es lo bastante rígida para romperse, así que en su interior
pueden empezar \omterm{def:g4:volcanoes-earthquakes-first-look:earthquake}{terremotos}, a cientos de kilómetros de profundidad.
\end{solution}

\begin{solution}{exo:g7:earthquakes:10}
(a) $2500 \div 800 = 3.125$~horas: 3~horas y $0.125 \times 60 = 7.5$
minutos, unas 3~horas y 8~minutos. (b) 3~h~8~min $-$ 2~h~17~min $=$
51~minutos más tarde. (c) La ola tardó horas en llegar, mientras que un
gran \omterm{def:g4:volcanoes-earthquakes-first-look:earthquake}{terremoto} submarino se localiza en pocos minutos: un aviso habría
podido llevar a la gente a lugares altos mucho antes de que llegara la ola.
\end{solution}

\begin{solution}{exo:g7:earthquakes:11}
(a) $400 \div 800 = 0.5$~horas: 30~minutos. (b) $20 \div 40 = 0.5$~horas:
30~minutos. (c) 60~minutos. (d) Los últimos \qty{20}{km} duran tanto como
los primeros \qty{400}{km}: la ola frena muchísimo en las aguas poco
profundas. (e) El frente de la ola frena mientras el agua que viene detrás
sigue llegando deprisa: el agua se amontona, y la ola crece en altura.
\end{solution}

\begin{solution}{exo:g7:earthquakes:12}
(a) Tres pasos: $32 \times 32 \times 32 = 32\,768$, unas
\textbf{33\,000}~veces. (b) $32\,768 \div 127 \approx 258$: unos
260~años de todos los \omterm{def:g4:volcanoes-earthquakes-first-look:earthquake}{terremotos} de \omterm{def:g7:earthquakes:magnitude}{magnitud}~6 del mundo. (c) Casi toda la
energía la liberan los pocos \omterm{def:g4:volcanoes-earthquakes-first-look:earthquake}{terremotos} más grandes: un solo \omterm{def:g4:volcanoes-earthquakes-first-look:earthquake}{terremoto} de
\omterm{def:g7:earthquakes:magnitude}{magnitud}~9 pesa más que siglos de \omterm{def:g4:volcanoes-earthquakes-first-look:earthquake}{terremotos} de \omterm{def:g7:earthquakes:magnitude}{magnitud}~6.
\end{solution}

\begin{solution}{pb:g7:earthquakes:1}
\textbf{1.} A: $45 - 25 = 20$~s. B: $29 - 16 = 13$~s. C: $34 - 19 =
15$~s.

\textbf{2.} La estación B: su retraso es el más corto, y cuanto más corto
es el retraso, más cerca está el \omterm{def:g4:volcanoes-earthquakes-first-look:earthquake}{terremoto}. (Sus ondas más rápidas también
llegaron primero.)

\textbf{3.} A: $8 \times 20 = \qty{160}{km}$. B: $8 \times 13 =
\qty{104}{km}$. C: $8 \times 15 = \qty{120}{km}$.

\textbf{4.} Todos los puntos situados a \qty{160}{km} de A: un círculo de
radio \qty{160}{km} centrado en A. El \omterm{def:g7:earthquakes:focus}{epicentro} está en algún punto de él.

\textbf{5.} A: $160 \div 20 = 8$~cuadrados. B: $104 \div 20 = 5.2$~cuadrados.
C: $120 \div 20 = 6$~cuadrados.

\textbf{6.} E está $3 + 6 = 9$~cuadrados al este y 14.5~cuadrados al norte
de O, en el \omterm{def:g2:water-on-earth:ocean}{mar}, entre las tres estaciones; los círculos de A y de B pasan
por él.

\textbf{7.} Bajo el \omterm{def:g2:water-on-earth:ocean}{mar}. Un movimiento brusco del fondo del \omterm{def:g2:water-on-earth:ocean}{mar} puede
provocar un \omterm{def:g7:earthquakes:tsunami}{tsunami}: hay que avisar a las \omterm{def:g1:landscapes-around-us:water}{costas} y a la isla.

\textbf{8.} Un paso de \omterm{def:g7:earthquakes:magnitude}{magnitud}: unas 32~veces más energía.

\textbf{9.} $12 \times 20 = \qty{240}{km}$.

\textbf{10.} $240 \div 800 = 0.3$~horas, y $0.3 \times 60 = 18$~minutos.

\textbf{11.} 5~minutos después del \omterm{def:g4:volcanoes-earthquakes-first-look:earthquake}{terremoto}.

\textbf{12.} $18 - 5 = 13$: los habitantes de Port Hale tienen unos
\textbf{13~minutos} para llegar a un lugar alto. (La ola frena en las aguas
poco profundas cerca de la isla, lo que da un poco más de tiempo, pero
nadie debe contar con ello: se marchan enseguida.)
\end{solution}
